\subsection{Review of Fourier analysis}
The fourier transform well-defined is 
\begin{align*}
    \hat{f}(a) := \frac{1}{\sqrt{2\pi}} \int_{-\infty}^\infty e^{-i a x}f(x)dx
\end{align*}
Sometimes $F[f]$ instead of $\hat{f}$ \\
\textbf{Inverse Fourier transform : } $\overset{u}{f}$ well-defined is 
\begin{align*}
    \overset{u}{f}(x) =  \frac{1}{\sqrt{2\pi}} \int_{-\infty}^\infty  e^{iax}f(a)da
\end{align*}
Sometimes $F^{-1}[f]$ instead. 
\subsubsection{Properties of Fourier transform}
\begin{itemize}
    \item Linearity : $F[\lambda f + \mu g] = \lambda F[f] + \mu F[g]$, whenever $\lambda, \mu \in \mathbb{R}$
    \item Time shift: Let $g(x) = f(x-x_0)$ then 
    \begin{align*}
        \hat{g}(a) = e^{-iax_0}\frac{1}{\sqrt{2 \pi}}\int_{-\infty}^\infty e^{-iay}f(y)dy = e^{-iax_0}\hat{f}(a)
    \end{align*}
    \item Frequency shift : Let $g(x) = e^{ia_0x}f(x)$ for $a_0 \in \mathbb{R}$ 
    \begin{align*}
        \hat{g}(a) = \frac{1}{\sqrt{2 \pi}}\int_{-\infty}^\infty  e^{-(a-a_0)x}f(x)dx = \hat{f}(a-a_0) 
    \end{align*}
    \item rescaling/dilation : Let $g(x) = f(\lambda x)$ where $\lambda \in \mathbb{R}\setminus \{0\}$
    \begin{align*}
        \hat{g}(a) = \frac{1}{\sqrt{2 \pi}}\int_{-\infty}^\infty e^{-iax}f(\lambda x)dx
    \end{align*} 
    to recover $\hat{f}$ want to substitute $y = \lambda x$. Two cases
    \begin{itemize}
        \item if $\lambda > 0$, we got $\int_{-\infty}^\infty \cdots \frac{1}{\lambda}dy$
        \item if $\lambda < 0$, we got $\int_{\infty}^{-\infty} \cdots \frac{1}{\lambda}dy = - \int_{-\infty}^\infty \cdots \frac{1}{\lambda}dy = \int_{-\infty}^\infty \cdots \frac{1}{|\lambda |}dy$
    \end{itemize}
    In both cases : 
    \begin{align*}
        \hat{g}(a) = \frac{1}{\sqrt{2 \pi}}\int_{-\infty}^\infty e^{-iay/\lambda}f(y)\frac{1}{|\lambda|}dy = \frac{1}{|\lambda|} \hat{f}(\frac{a}{\lambda})
    \end{align*}
    \item derivatives  : Let $g(x) = f'(x)$. Then 
    \begin{align*}
        \hat{g}(a) = ia\hat{f}(a)
    \end{align*}
    More generally
    \begin{align*}
        F[f^{(u)}](a) = (ia)^u F[f](a)
    \end{align*}
    So then Fourier transform converts differential equation with constant coefficients into algebric equations on frequency domain. 
\end{itemize}
\subsubsection{Convolution}
Given $f, g : \mathbb{R} \rightarrow \mathbb{C},$ the convolution $f \ast g$, whenever well-defined is 
\begin{align*}
    f \ast g (x) = \int_{-\infty}^\infty f(y)g(x-y) dy
\end{align*}
Some properties of convolution:
\begin{itemize}
    \item $f \ast g = g \ast f$
    \item $f \ast (g \ast h) = (f \ast g) \ast h$
    \item $f \ast (g+h) = f\ast g + f \ast h$
    \item $(f \ast g)' = f' \ast g = f \ast g'$
\end{itemize}
\begin{itemize}
    \item Convolution : 
    \begin{align*}
    F[f \ast g] = \sqrt{2 \pi} F[f]F[g]
    \end{align*}
    \item Gaussian function : if $f(x) = e^{-x^2/2},$ 
    \begin{align*}
        \hat{f}(a) = e^{-a^2/2}
    \end{align*}
    \item Plancherel's theorem : 
    \begin{align*}
        \int_{-\infty}^\infty |f(x)|^2dx = \int_{-\infty}^\infty  |\hat{f}(a)|^2da
    \end{align*}
\end{itemize}

\subsection{Laplace analysis}
Fourier transform computed for signals on $\mathbb{R}$ \\
Laplace transform: Computed for signals defined on $\mathbb{R_+}$. \\
Both transform differential equ. into algebric equ. \\
Let $f: \mathbb{R_+} \rightarrow \mathbb{C}$, where well-defined, the \textbf{Laplace transform} of $f$ at $z \in \mathbb{C}$ is 
\begin{align*}
    \mathcal{L}[f](z) := \int_0^\infty f(t) e^{-zt}dt
\end{align*}
Sometimes Laplace transform of functions are written by capital letters. 
\subsubsection{Convergence criteria}
Of course the above integral well-edfined if $\int_{-\infty}^\infty |f(t)||e^{-zt}dt| < \infty$ Assume that is $\gamma_0 \in \mathbb{R}$ st
\begin{align*}
    \int_0^\infty |f(t)|e^{-\gamma_0t}dt < \infty 
\end{align*}
(This means $|f(t)| < e^{\gamma_0 t}$ as $t \rightarrow \infty)$
\begin{align*}
    |\mathcal{L}[f](z)| \leq \int_0^\infty |f(t)||e^{-zt}|dt \\
    = \int_0^\infty |f(t)||e^{-Re(z) t}|dt \\
    \leq \int_0^\infty |f(t)|e^{-\gamma_0t}dt < \infty 
\end{align*}
Thus if such a $\gamma_0$ exists, $\mathcal{L}[f](z)$ is defined at least for $z \in \{Re > \gamma_0\}$
\subsubsection{Definition} A $\gamma_0 \in \mathbb{R}$ such that $\mathcal{L}[f](z)$ is well-defined at least on $\{Re > \gamma_0\}$ is called \textbf{abscission of convergence} \\
Let $\gamma_0$ abscisse of convergence. Thus $\mathcal{L}[f]$ is holomorphic on $\{Re > \gamma_0\}$. Set $F(z) = \mathcal{L}[f](z)$. Thus 
\begin{align*}
    \frac{d}{dz}F(z) = \frac{d}{dz}\int_0^\infty f(z)e^{-zt}dt \\
    = \int_0^\infty f(t) (-t)e^{-zt}dt = - \mathcal{L}[tf(t)](z)
\end{align*}
In short 
\begin{align*}
    \mathcal{L}[f(t)]' = - \mathcal{L}[tf(t)]
\end{align*}